\subsection{Parabola: equal distance from a focus and a line}

Fix a focus $F=(p,0)$ and a directrix $D:x=-p$, where $p>0$. The parabola is the
set of points $P=(x,y)$ satisfying
\[
\boxed{d(P,F)=d(P,D)}.
\]
Since
\[
d(P,F)=\sqrt{(x-p)^2+y^2},
\qquad
d(P,D)=x+p,
\]
squaring the equality gives
\[
(x-p)^2+y^2=(x+p)^2,
\]
and therefore
\[
\boxed{y^2=4px}.
\]

\begin{center}
\begin{tikzpicture}[scale=0.88,>=stealth]
    \draw[axis,->] (-2.5,0)--(6.3,0) node[right] {$x$};
    \draw[axis,->] (0,-4.0)--(0,4.2) node[above] {$y$};
    \draw[gray,thick] (-1.2,-3.6)--(-1.2,3.6)
        node[above] {$D$};
    \draw[subset,domain=-3.3:3.3,samples=120]
        plot ({0.42*\x*\x},{\x});
    \coordinate (F) at (1.2,0);
    \coordinate (P) at (2.7,2.55);
    \coordinate (H) at (-1.2,2.55);
    \draw[blue!70!black,line width=1.1pt] (F)--(P)
        node[midway,below right] {$d(P,F)$};
    \draw[orange!85!black,line width=1.1pt] (P)--(H)
        node[midway,above] {$d(P,D)$};
    \fill[geometricSubsetColor] (F) circle (2.2pt) node[below] {$F$};
    \fill[geometricSubsetColor] (P) circle (2.2pt) node[above right] {$P$};
    \fill[geometricSubsetColor] (0,0) circle (2.2pt) node[below left] {vertex};
\end{tikzpicture}
\end{center}

A convenient parameterisation is
\[
\boxed{x=pt^2,\qquad y=2pt}.
\]
Substitution gives $y^2=4p^2t^2=4px$, so every parameter value produces a point
on the curve.

\begin{interpretationbox}[Why the parabola is the boundary case]
The parabola is neither closed like an ellipse nor split into two branches like
a hyperbola. It will later appear exactly at the transition between bounded and
unbounded focal motion.
\end{interpretationbox}
